\documentclass[parskip=full]{scrartcl}
\usepackage[top=2.5cm, bottom=2.5cm, left=2.5cm, right=2.5cm]{geometry}
\usepackage[utf8]{inputenc}
\usepackage[T1]{fontenc}
\usepackage[german]{babel}
\usepackage{hyperref}
\usepackage[toc, nonumberlist]{glossaries}
\usepackage{graphicx}
\usepackage{enumitem}
\usepackage{float}
\usepackage{color}

\title{OSIP - Softwaredesign}
\subtitle{OPC UA Simulator for Industrial Plants}
\author{
    M. Armbruster\\
    D. Kahles\\
    H. Lehmann\\
    M. Schwarzmann\\
    N. Wilhelm
}

\begin{document}
\maketitle

\vspace{20px}
\begin{center}
  \includegraphics[scale=0.4]{../icon.png}
\end{center}
\pagebreak
\tableofcontents
\pagebreak

\section{Einleitung}
% Aufteilung in die Pakete erklären

\section{Klassenbeschreibungen}
% JavaDoc

\section{Abläufe}
% Sequenzdiagramme

\section{Änderungen zum Pflichtenheft}

\section{\"Ubernommene Kann-Kriterien}
Im folgenden werden die Kann-Kriterien aus dem Pflichtenheft noch einmal alle aufgelistet. Zu jedem der Kriterien wird
kurz zusammengefasst, was genau es bedeutet. Es werden erst die implementierten und anschlie{\ss}end die nicht implementierten
Kriterien gelistet.

%Implementiert:
%- Logging von Fehler bei der Ausführung stdout
%- Szenarien
%- Jitter
%- Zeitintervall
%- Alarme ein-/ausschalten (und das abspeichern)
%- Logging für Alarme in der Überwachungskonsole
\subsection{Implementierte Kann-Kriterien}
Die implementierten Kann-Kriterien sind wie folgt:

Sowohl in der \"Uberwachungskonsole als auch in der Fertigungssimulation werden Fehler bei der Programmausf\"uhrung nach
stdout geloggt.

In der Fertigungssimulation k\"onnen vorgefertigte Simulationsszenarien geladen werden. Diese sind in eigenen Dateien definiert.
Die Ausführung der Simulationsszenarien f\"uhrt dazu, dass die Steuervariablen sich nach Vorgaben des Szenarios ohne
weitere Aktion des Benutzers \"andern.

In den Zuflusswerten der Fertigungssimulation ist ein Jitter eingebaut. Dieser sorgt f\"ur leichte Variationen in den
Zuflussmengen und der Temperatur der Zufl\"usse. Das Ausma{\ss} des Jitters ist nicht einstellbar.

Innerhalb der \"Uberwachungskonsole kann der Benutzer selbst einstellen, in welchem Zeitintervall die Werte der
Fertigungssimulation abgefragt werden. Alle Werte werden im jeweils gleichen Intervall abgefragt.

Die Alarme in der \"Uberwachungskonsole k\"onnen vom Benutzer jeweils einzeln zu- und abgeschaltet werden. Die Menge der
aktivierten Alarme wird bei Beenden der \"Uberwachungskonsole persistiert. Beim n\"achsten Start der \"Uberwachungskonsole
werden die gleichen Alarme automatisch wieder aktiviert.

Die Alarme, die von der \"Uberwachungskonsole empfangen werden, werden in einer Textausgabe innerhalb der \"Uberwachungskonsole
geloggt.

%nicht implementiert:
%- andere Sprache
%- Alarmschwellen konfigurieren
%- Jitter einstellen
\subsection{Nicht implementierte Kann-Kriterien}
Die nicht implementierten Kann-Kriterien sind wie folgt:

%Sollen wir es moeglich machen, eigene Sprachdateien einzulesen?
Es stehen keine verschiedenen Sprachoptionen f\"ur \"Uberwachungskonsole und Fertigungssimulation zur Verf\"ugung. Die einzige
vorhandene Sprache ist Englisch.

Die vorgegebenen Alarme k\"onnen nicht ver\"andert werden. Insbesondere ist es nicht m\"oglich, eigene Alarme hinzuzuf\"ugen
oder die Schwellenwerte der vorgegebenen Alarme zu modifizieren.

Der Jitter in den Simulationswerten ist fest vorgegeben und kann in seiner Intensit\"at nicht ver\"andert werden.

\section{Klassendiagramme}
Im Anhang.

\section{Verwendete Entwurfsmuster}
% Aufzählen: Geheimnisprinzip, MVC etc.

\end{document}
